% D04-009. AI-01 P0197--P0214, сохранённый отчёт и код проверки.
\section{Анализ обученных физических моделей}
\label{sec:ch4_physical_validation}
\conceptsource{Статья 1. Обученные функции, отключение арифметических ветвей и обнуление связей раскрывают использование физических преобразований. Метод проверки уточнён по сохранённому опыту.}

% D04-011 R4-08 BEGIN
\ifdefined\AuthorReview
\par\smallskip\begingroup\small\raggedright\color{red!60!black}
\noindent\textbf{Примечание рецензента [R4-08], важно.} После подробного описания восьми ранних архитектур читателю не хватает сопоставимого результата их сравнения. Отключение ветвей у одной PhysicsKAN показывает использование этих ветвей, но не превосходство семейства над базовыми сетями. Радар и инженерные столбцы AI-01 расходятся сильнее округления; их нельзя объединять. Перед итогами главы либо восстановить одну согласованную серию, либо явно ограничить вклад этого блока разработкой и исследованием устройства, без рейтинга архитектур.\par\endgroup\smallskip
\fi
% D04-011 R4-08 END
\ifdefined\AuthorReview
\par\smallskip\begingroup\small\raggedright\color{blue!55!black}
\noindent\textbf{Ответ на R4-08, D04-012-B.} Согласована связь каждого архитектурного приёма с назначением и доступной проверкой. Итог ссылается на отключение ветвей, сравнение входов и прикладную серию ОЗЗ. Расходящиеся сводные графики не объединены в общий численный рейтинг.\par\endgroup\smallskip
\fi

(Примечание автора: Со мной согласовано, ожидается проверка рецензента.)


Структура физических блоков задаёт определённый способ
обработки входных величин. После обучения возникает другой
вопрос, насколько модель использует эти преобразования
и какие зависимости формирует в последующих слоях.
В исследовании архитектур\footnote{Евдаков А.Е.
Архитектурная модернизация интеллектуальных алгоритмов
релейной защиты. Внедрение законов электротехники
в граф вычислений.} рассмотрены обученные функции KAN,
отключение групп признаков и обнуление слабых связей.
Эти проверки дополняют общие показатели распознавания.

\subsection{Обученные функции на связях}
\label{sec:ch4_learned_functions}
\conceptsource{Статья 1, рисунок 11. Четыре кривые извлечены из реального рисунка обученной PhysicsKAN; форма не синтезировалась, масштаб каждого поля свой.}


В KAN отклик одной связи задаётся отдельной одномерной
функцией. Поэтому после обучения можно непосредственно
рассмотреть её форму, а не только итоговую оценку класса.
На рисунке~\ref{fig:ch4_learned_kan_functions} показаны
выбранные функции обученной PhysicsKAN.
Это четыре функции из исходного рисунка обученной модели.
Их форма сохранена, а каждое поле имеет отдельный масштаб.
Кривые показывают реальные отклики связей, полученные
после обучения, а не условные примеры желаемой характеристики.
У них видны участки разной крутизны, смены знака
и немонотонного изменения отклика.

\iffalse
(Примечание автора: Это же не реально обученные, да? Но выглядят красиво подставные конечно. Можно подумать, что оставить, ибо в работе вроде бы реальные были, но они похуже.) **[AA095: выбор рисунка принят автором 27.10; история сохранена]**
\fi
\authorreviewnote{AA095}

\begin{figure}[htbp]
 \centering
 \includegraphics[width=\textwidth]{ch04_learned_kan_functions.pdf}
 \caption{Формы отдельных обученных функций на внутренних
 связях PhysicsKAN. Масштаб каждого поля выбран отдельно}
 \label{fig:ch4_learned_kan_functions}
\end{figure}
\figurereview{По уточнению автора 27.10 сохраняется компактная перерисовка реальных кривых. Источник рисунок 11 статьи 1, поля 1, 2, 5 и 9 из fig11\_kan\_splines\_article.svg. Перенесены исходные векторные пути, без генерации формы и без повторного инференса. Автоматические названия типов кривых исходника не использованы.}


Участок с большим изменением отклика при небольшом изменении
входа может участвовать в формировании границы решения.
Более пологий участок задаёт меньшую чувствительность
к соответствующему признаку.
Суммирование таких откликов позволяет получить составную
нелинейную характеристику, которая затем используется
следующими слоями и выходным классификатором.

Физическая интерпретация зависит от положения связи.
На входе сети величина может сохранять связь с измеренным
током, напряжением или явно рассчитанным признаком.
Во внутренних слоях она обычно является результатом
нескольких преобразований и смешивания каналов.
Связь выбранной кривой с конкретной уставкой
или электрической величиной устанавливается
по полному пути её формирования и проверке режимов.

Отдельная форма кривой также не определяет её вклад
во все решения модели. Существенны величина отклика,
диапазон входа на реальных примерах и компенсация
другими связями. Поэтому визуальное рассмотрение
функций дополняется проверками с изменением состава
обработки и оценкой результата на осциллограммах.

\subsection{Отключение исходных и производных признаков}
\label{sec:ch4_physical_ablation}
\conceptsource{Статья 1. Отключение групп признаков у фиксированной обученной модели, без нового обучения каждой конфигурации.}


Для обученной PhysicsKAN средней сложности
рассмотрено несколько вариантов входного физического блока.
В полной конфигурации сохраняются исходные координаты,
их попарные произведения и отношения.
Затем отдельно обнуляются отклики ветвей произведений,
отношений, обеих арифметических ветвей либо исходного пути.
Остальные коэффициенты модели остаются фиксированными.
После каждого изменения выполняется оценка на той же
проверочной выборке, без повторного обучения.

Результаты представлены на
рисунке~\ref{fig:ch4_physical_feature_ablation}.
На исходном рисунке обозначения $S$ и $Y$ соответствуют
ветвям произведений и отношений входных координат.
Их физическая трактовка уточняется ниже.
Полная конфигурация имеет Macro-F1 0,602.
При отключении произведений показатель уменьшается до 0,397,
при отключении отношений до 0,593.
Одновременное отключение обеих арифметических ветвей
даёт 0,396. Сохранение только производных признаков
без исходного пути даёт 0,336.

\begin{figure}[htbp]
 \centering
 \includegraphics[width=\textwidth]{ch04_feature_ablation_article.png}
 \caption{Изменение Macro-F1 при отключении групп признаков
 входного блока обученной PhysicsKAN}
 \label{fig:ch4_physical_feature_ablation}
\end{figure}
\figurereview{Возвращён исходный рисунок 13 статьи 1. На нём S и Y условно обозначают ветви произведений и отношений. Перед выпуском заменить S=U*I и Y=I/U на названия ветвей, поскольку для текущих вещественных координат это не буквальные комплексные мощность/проводимость.}


В этой модели произведения оказывают заметное влияние
на сформированное решение. Исходные признаки также
сохраняют самостоятельную роль, поскольку их отключение
приводит к наибольшему снижению показателя.
Результат поддерживает совместную передачу
исходных и производных величин в рассмотренном блоке.

Относительно небольшое изменение при отключении отношений
характеризует использование этой ветви обученной моделью.
Оно само по себе не определяет качество новой архитектуры,
изначально обученной без отношений.
После удаления части входа распределение скрытых признаков
меняется, а коэффициенты не успевают приспособиться.
Проверка с отдельным обучением каждой конфигурации
отвечает на другой вопрос и рассматривается отдельно.

В данном вещественном блоке выполняются произведения
и отношения входных координат.
Их электротехническая трактовка зависит от представления
сигнала. При полярном входе произведение двух отдельных
координат не равно автоматически комплексной мощности.
Комплексная арифметика выделена в самостоятельные
архитектурные варианты из
подраздела~\ref{sec:ch4_complex_ops}.

\subsection{Обнуление слабых связей}
\label{sec:ch4_kan_pruning}
\conceptsource{Статья 1 и сохранённый отчёт. Зависимость качества от обнулённых связей; уменьшение времени отдельно не измерялось.}


Отклики отдельных связей KAN имеют разный масштаб
на выбранных примерах. Это позволяет оценивать,
какие связи можно обнулить с небольшим изменением
результата. В рассмотренном опыте величина связи
оценивалась по среднему модулю её отклика.
Затем к копии обученной модели применялся выбранный порог,
а качество повторно оценивалось без дообучения.

На рисунке~\ref{fig:ch4_kan_pruning} показана зависимость
Macro-F1 от доли сохранённых связей.
При сохранении 93,7\% связей получена оценка 0,606
при исходной 0,602. Различие мало и характеризует
рассмотренную проверочную выборку.
При сохранении 76,9\% связей показатель составляет 0,582,
а при сохранении 59,8\% уменьшается до 0,550.
Более сильное обнуление приводит к значительной
потере качества.

\begin{figure}[htbp]
 \centering
 \includegraphics[width=\textwidth]{ch04_kan_pruning_report.png}
 \caption{Качество обученной PhysicsKAN при обнулении связей
 с малым средним откликом}
 \label{fig:ch4_kan_pruning}
\end{figure}
\figurereview{Исходная кривая сохранённого отчёта того же опыта. Перевести Pruning Curve и F1 (macro). Ось показывает сохранённые связи, в отличие от прежней перерисовки по доле обнулённых.}


Таким образом, запас для сокращения модели определяется
совместно по числу сохраняемых связей и ошибкам распознавания.
Слабый отклик на одних примерах может быть существенным
для редкого повреждения. Поэтому средняя оценка качества
рассматривается вместе с отдельными классами и режимами.

Обнуление коэффициентов и уменьшение времени расчёта
также различаются. При сохранении прежних плотных матричных
операций нулевые коэффициенты могут продолжать участвовать
в тех же вычислениях. Практическое ускорение требует
подходящей реализации сокращённой структуры
и отдельного измерения на выбранной аппаратуре.

\subsection{Значение проверок для измерительных органов}
\label{sec:ch4_physical_checks_meaning}
\conceptsource{Итог опыта статьи 1. Разделены использование физических ветвей, возможность сокращения и требования к прикладной проверке.}


\iffalse
(Примечание автора: Если это краткий вывод – то явно надо переписать, как-то он мало понятен и слишком заумный, вероятно касается и других выводах в подглавах других, а не только здесь.) **[AA096: рассмотрено 27.10; ожидает проверки]**
\fi
\authorreviewnote{AA096}

Проверки показывают, что физические ветви действительно
используются обученной моделью. В рассмотренном опыте
наибольшее снижение качества связано с отключением
произведений и исходных признаков. Одних производных
величин для прежнего качества недостаточно.

Обнуление небольшого числа слабых связей почти
не меняет общий результат, а более сильное сокращение
увеличивает ошибки. Поэтому сохраняется совместная
передача исходных и рассчитанных признаков,
а сокращение сети оценивается по качеству распознавания.

Эти результаты относятся к составу обработки.
Далее рассматриваются сами переходные процессы,
на которых видны пропуски, лишние сигналы
и несовпадения границ событий.
